\documentclass{article}

\usepackage{float}
\restylefloat{table}

\usepackage{booktabs}

\title{Team Contributions: Rev 0\\\progname}

\author{\authname}

\date{}

\input{../Comments}
\input{../Common}

\begin{document}

\maketitle

This document summarizes the contributions of each team member for the Rev 0
Demo.  The time period of interest is the time between the PoC demo and the Rev
0 demo; the contributions prior to the PoC are NOT included.

\section{Demo Plans}

%\wss{What will you be demonstrating}
The demonstration for Rev 0 will focus on the core functionality of the Phyio Companion application, mainly the ability to record performed exercises, view past exercise data, and receive exercise recommendations from the patient's assigned therapist. 

The physiotherapist interface will also be demonstrated, showcasing how therapists can assign exercises to patients and monitor their progress through the application, including leaving comments on exercise logs and adjusting exercise plans as needed.

The demo will also include a brief overview of the application's architecture and design choices made during development.

\section{Team Meeting Attendance}

%\wss{For each team member how many team meetings have they attended over the
%time period of interest.  This number should be determined from the meeting
%issues in the team's repo.  The first entry in the table should be the total
%number of team meetings held by the team.}

\begin{table}[H]
\centering
\begin{tabular}{ll}
\toprule
\textbf{Student} & \textbf{Meetings}\\
\midrule
Total & 6\\
Matthew & 6\\
Cieran & 5\\
Vaisnavi & 6\\
Maham & 5\\
Eman & 5\\
\bottomrule
\end{tabular}
\end{table}

%\wss{If needed, an explanation for the counts can be provided here.}
The additional attendance from Matthew and Vaisnavi is due to a single coding sub-team meeting.

\section{Supervisor/Stakeholder Meeting Attendance}

%\wss{For each team member how many supervisor/stakeholder team meetings have
%they attended over the time period of interest.  This number should be determined
%from the supervisor meeting issues in the team's repo.  The first entry in the
%table should be the total number of supervisor and team meetings held by the
%team.  If there is no supervisor, there will usually be meetings with
%stakeholders (potential users) that can serve a similar purpose.}

\noindent \textbf{Supervisor's Name: Kevin Yu} 
\begin{table}[H]
\centering
\begin{tabular}{ll}
\toprule
\textbf{Student} & \textbf{Meetings}\\
\midrule
Total & 0\\
Matthew & 0\\
Cieran & 0\\
Vaisnavi & 0\\
Maham & 0\\
Eman & 0\\
\bottomrule
\end{tabular}
\end{table}

%\wss{If needed, an explanation for the counts can be provided here.}
Supervisor information provided before the PoC demo seemed sufficient for an initial implementation. There will be future meetings to ensure the application is meeting stakeholder requirements and expectations, however this period of time has been spent with existing information.
\section{Lecture Attendance}

% \wss{For each team member how many lectures have they attended over the time
% period of interest.  This number should be determined from the lecture issues in
% the team's repo. You can find the number of lectures in the time period of
% interest by looking at the
% \href{https://calendar.google.com/calendar/u/0/embed?src=rnboqiaki1k2la7rpu3bn0um58@group.calendar.google.com&ctz=America/Toronto}
% {Google calendar} for the capstone course.}

%\wss{NOTE: There will be approximately 1 lecture between the POC and Rev0 demos}

\begin{table}[H]
\centering
\begin{tabular}{ll}
\toprule
\textbf{Student} & \textbf{Lectures}\\
\midrule
Total & 1\\
Matthew & 1\\
Cieran & 1\\
Vaisnavi & 1\\
Maham & 1\\
Eman & 1\\
\bottomrule
\end{tabular}
\end{table}

%\wss{If needed, an explanation for the lecture attendance can be provided here.}

\section{TA Document Discussion Attendance}

% \wss{For each team member how many of the informal document discussion meetings
% with the TA were attended over the time period of interest.}

\noindent \textbf{TA's Name: Rashad Bhuiyan}

\begin{table}[H]
\centering
\begin{tabular}{ll}
\toprule
\textbf{Student} & \textbf{Lectures}\\
\midrule
Total & 1\\
Matthew & 0\\
Cieran & 1\\
Vaisnavi & 1\\
Maham & 1\\
Eman & 1\\
\bottomrule
\end{tabular}
\end{table}

This meeting was scheduled outside of the usual time due to a snow day, and the team could not find a time with no scheduling conflicts before the due date of Rev 0 design documentation. Therefore, this meeting was scheduled outside of Matthew's availability to allow all of the documentation subteam to attend.

\section{Commits}

% \wss{For each team member how many commits to the main branch have been made
% over the time period of interest.  The total is the total number of commits for
% the entire team since the beginning of the term.  The percentage is the
% percentage of the total commits made by each team member.}

\begin{table}[H]
\centering
\begin{tabular}{lll}
\toprule
\textbf{Student} & \textbf{Commits} & \textbf{Percent}\\
\midrule
Total & 81 & 100\% \\
Matthew & 0 & 0\% \\
Cieran & 16 & 19.8\% \\
Vaisnavi & 29 & 35.8\% \\
Maham & 36 & 44.4\% \\
Eman & 0 & 0\% \\
\bottomrule
\end{tabular}
\end{table}

% \wss{If needed, an explanation for the counts can be provided here.  For
% instance, if a team member has more commits to unmerged branches, these numbers
% can be provided here.  If multiple people contribute to a commit, git allows for
% multi-author commits.}

Eman's commitment to the documentation were made through Cieran, who forgot to co-author his commits. Matthew's section of the code was not included at the time of this report, but will be added and spoken on during the demonstration.

\section{Issue Tracker}

% \wss{For each team member how many issues have they authored (including open and
% closed issues (O+C)) and how many have they been assigned (only counting closed
% issues (C only)) over the time period of interest.}

\begin{table}[H]
\centering
\begin{tabular}{lll}
\toprule
\textbf{Student} & \textbf{Authored (O+C)} & \textbf{Assigned (C only)}\\
\midrule
Matthew & 0 & 0 \\
Cieran & 0 & 0 \\
Vaisnavi & 0 & 0 \\
Maham & 0 & 0 \\
Eman & 0 & 0 \\
\bottomrule
\end{tabular}
\end{table}

%\wss{If needed, an explanation for the counts can be provided here.}
Over this period of development, the team focused on previously assigned tasks and sections of the project and was not tracked through issues.

\section{CICD}

%\wss{Say how CICD is used in your project}
CICD is still being used to track compilation of all .tex files and has unit testing setups ready for between Rev 0 and Rev 1. Currently however, the CICD remains focused on .tex files.

There was a change made to disregard .pdf files from being tracked in the template repository, however our team has not adopted those changes.

\section{Team Charter Trigger Items}

\begin{itemize}
    \item Missing 3 consecutive meetings when deliverables are not the main discussion topic triggers a team meeting
    \item Missing 2 consecutive meetings when deliverables are the main topic triggers a team meeting
    \item Deliverables must be finalized 2 days prior to the deadline (internal deadline)
\end{itemize}
 
\subsection{Violations}
No violations have been recorded at this time.
 
\subsection{Plan to Address Violations}
If triggers are violated, the team will:

\begin{enumerate}
    \item Hold a team meeting to discuss the attendance or deadline issues.
    \item Review the circumstances with the member involved.
    \item For extended emergencies ($>$2 weeks), bring the matter to teaching staff to ensure workload remains manageable for the team.
    \item If triggers are found to be too weak, strong, or ambiguous, they will be revised through the team's decision-making process (this will be done through a unanimous vote or majority vote with all members heard).
\end{enumerate}

\section{Additional Productivity Metrics}

\textit{Duration between pull requests:} Over the course of the project, our team learned that more frequent pull requests would be beneficial to the workflow. Frequent merges will ensure minimal merge conflicts towards the end of the deliverable and guarantee that any output is up-to-date. Merge conflicts are time-consuming and troublesome to handle, and implementing this strategy will result in an overall improvement in the quality of our work. 


\end{document}